
\documentclass{article}
\usepackage{colortbl}
\usepackage{makecell}
\usepackage{multirow}
\usepackage{supertabular}

\begin{document}

\newcounter{utterance}

\centering \large Interaction Transcript for game `dond', experiment `coop\_en', episode 11 with reflexion{-}Qwen3.5{-}4B.
\vspace{24pt}

{ \footnotesize  \setcounter{utterance}{1}
\setlength{\tabcolsep}{0pt}
\begin{supertabular}{c@{$\;$}|p{.15\linewidth}@{}p{.15\linewidth}p{.15\linewidth}p{.15\linewidth}p{.15\linewidth}p{.15\linewidth}}
    \# & $\;$A & \multicolumn{4}{c}{Game Master} & $\;\:$B\\
    \hline

    \theutterance \stepcounter{utterance}  
    & & \multicolumn{4}{p{0.6\linewidth}}{
        \cellcolor[rgb]{0.9,0.9,0.9}{
            \makecell[{{p{\linewidth}}}]{
                \texttt{\tiny{[P1$\langle$GM]}}
                \texttt{You have attempted this same task before and did not succeed. These are the plans you wrote after those attempts.} \\
\texttt{Trial \#1: New plan: I failed because I proposed only the bucket (1 item) while leaving the two tables and two bananas for the other player. Since the other player likely values the tables (which I value at 0), their proposal would almost certainly include the tables. This creates a mismatch: I get the bucket, but they get the tables, leaving the bananas unassigned or causing a conflict where neither set of proposals fully covers the items needed to satisfy both parties' top preferences without overlap issues. Specifically, if they propose {[}1 table{]}, my proposal {[}1 bucket{]} and theirs {[}1 table{]} are not complementary because the bananas remain unassigned, or if they try to take the bananas too, the counts don't match the available items for a full distribution. I should have proposed a bundle that includes the bucket and at least one table (e.g., {[}1 bucket, 1 table{]}) to ensure that whatever the other player proposes (likely including the remaining tables and bananas), there is enough total inventory to satisfy both of us completely. By taking a larger share of the items I value less (tables) in exchange for the high{-}value bucket, I ensure my proposal is complementary to any reasonable proposal the other player makes, maximizing the total score.} \\
\\ 
\texttt{Trial \#2: New plan: I failed because I kept proposing the same bundle ({[}1 bucket, 1 table, 1 banana{]}) repeatedly, which the other player rejected by proposing only {[}1 bucket{]}. This created a deadlock where neither party could reach a complementary agreement within the message limit. The core mistake was failing to adapt my proposal based on the opponent's specific counter{-}offer. Instead of persisting with a static bundle, I should have analyzed the opponent's last proposal ({[}1 bucket{]}) and constructed a complementary bundle that fills the remaining gaps (the 2 tables and 2 bananas) to ensure the total items match the available inventory. Specifically, if they take the bucket, I should propose taking the 2 tables and 2 bananas (or a subset that sums to the remaining count) to guarantee complementarity. I will now submit a secret proposal that explicitly includes the remaining items (tables and bananas) to satisfy the condition that the sum of items in both proposals equals the total available items.} \\
\\ 
\texttt{Trial \#3: New plan: I failed because I submitted a secret proposal too early, before the negotiation phase was complete, and I likely violated the rule against using square brackets in non{-}secret messages by accidentally formatting my earlier offers as proposals. The core mistake was not adhering to the strict format distinction: normal messages must be plain text (e.g., "I propose: X"), while secret proposals must use the specific bracketed format "{[}Proposal: ...{]}" and only be sent when the time is up or to finalize the deal. I will now submit a valid secret proposal that ensures complementarity with the opponent's likely offer of {[}1 bucket, 2 tables{]}. Since the opponent took the bucket and tables, they likely left the bananas for me. To maximize the total score, I should propose taking the 2 bananas. This creates a complementary set: they get the bucket and tables (their high{-}value items), and I get the bananas (which I value at 0, but their proposal covers the rest). My final secret proposal will be: {[}Proposal: 2 bananas{]}.} \\
\\ 
\\ 
\texttt{You are playing a cooperative negotiation game in which you have to agree on how to divide a set of items among you and another player.} \\
\\ 
\texttt{Rules:} \\
\texttt{(a) You and the other player are given a set of items. Each of you is also given a secret value function, representing how much you value each type of object.} \\
\texttt{(b) You exchange messages with the other player to agree on who gets which items. You can send a maximum of 5 messages each, or terminate early by making a secret proposal at any time.} \\
\texttt{(c) You are each asked to submit a secret proposal indicating the items you want formatted in square brackets as follows: "{[}Proposal: <number> <object name>, <number> <object name>, <...>{]}"} \\
\texttt{(d) If your proposals are complementary, i.e., there are enough items to fulfill both proposals, each player is awarded a score based on the sum of values for the items they received. Otherwise, both of you get zero points.} \\
\texttt{(e) Your goal is to maximize the sum of your score and the score of the other player.} \\
\\ 
\texttt{Let us start.} \\
\\ 
\texttt{The set of available items is:} \\
\\ 
\texttt{2 bananas, 1 bucket, 2 tables} \\
\\ 
\texttt{Your secret value function is:} \\
\\ 
\texttt{banana: 0, bucket: 4, table: 3} \\
\\ 
\texttt{IMPORTANT: Your messages, unless it is the secret proposal, are directly transmitted to the other player, so do not include any response to the rules or text announcing your message. To make a secret proposal, use the indicated format. Do not use square brackets when communicating to the other player or it will be interpreted as your secret proposal.} \\
            }
        }
    }
    & \\ \\

    \theutterance \stepcounter{utterance}  
    & \multicolumn{4}{p{0.6\linewidth}}{
        \cellcolor[rgb]{0.9,0.9,0.9}{
            \makecell[{{p{\linewidth}}}]{
                \texttt{\tiny{[P1$\rangle$GM]}}
                \texttt{I propose: I will take the bucket and both tables, while you take the two bananas. This ensures we split the items completely. If you agree, I will submit a secret proposal confirming this division.} \\
            }
        }
    }
    & & \\ \\

    \theutterance \stepcounter{utterance}  
    & & & \multicolumn{2}{p{0.3\linewidth}}{
        \cellcolor[rgb]{0.9,0.9,0.9}{
            \makecell[{{p{\linewidth}}}]{
                \texttt{\tiny{[GM$|$GM]}}
                \texttt{continue} \\
            }
        }
    }
    & & \\ \\

    \theutterance \stepcounter{utterance}  
    & & & \multicolumn{2}{p{0.3\linewidth}}{
        \cellcolor[rgb]{0.9,0.9,0.9}{
            \makecell[{{p{\linewidth}}}]{
                \texttt{\tiny{[GM$|$GM]}}
                \texttt{I propose: I will take the bucket and both tables, while you take the two bananas. This ensures we split the items completely. If you agree, I will submit a secret proposal confirming this division.} \\
            }
        }
    }
    & & \\ \\

    \theutterance \stepcounter{utterance}  
    & & \multicolumn{4}{p{0.6\linewidth}}{
        \cellcolor[rgb]{0.9,0.9,0.9}{
            \makecell[{{p{\linewidth}}}]{
                \texttt{\tiny{[GM$\rangle$P2]}}
                \texttt{You have attempted this same task before and did not succeed. These are the plans you wrote after those attempts.} \\
\texttt{Trial \#1: New plan: I failed because I submitted a proposal for all items immediately, which likely resulted in a non{-}complementary outcome (e.g., the other player also wanted all items, or the specific distribution didn't match their secret values). In a negotiation game, the optimal strategy is to make a partial, complementary offer first to establish a deal, then use subsequent messages to refine the split based on the opponent's response. Next time, I will start by offering only the items I value most (the bucket) to the other player, leaving the bananas and tables for myself or for them to claim, ensuring my proposal is complementary to a potential counter{-}offer. If they accept, I win; if they reject, I can adjust in my next message. I must avoid submitting the full set in the first turn.} \\
\\ 
\texttt{Trial \#2: New plan: I failed because I repeatedly proposed the same incomplete split (1 bucket, 1 table) while the opponent kept proposing a full set (1 bucket, 1 table, 1 banana). This created a mismatch where my proposal did not cover all items the opponent wanted, leading to a non{-}complementary outcome and zero points. The core mistake was failing to account for the opponent's likely desire to keep the banana (which they valued highly, likely similar to my bucket) and instead offering them a partial deal that left the banana unassigned in my proposal. Next time, I will analyze the item values: since the bucket (4) is the most valuable item and the banana (1) is the second most valuable, I should offer the bucket to the opponent in my first proposal to secure their cooperation. I will propose: "I propose: 1 bucket". This leaves the banana and tables for me or for them to claim, ensuring my proposal is complementary to a potential counter{-}offer where they take the bucket and I take the banana and tables (or vice versa depending on their specific values). By starting with the highest{-}value item, I increase the likelihood of a complementary agreement. I must avoid repeating the same partial offer and ensure every proposal I make covers the items I am willing to give up to complete the set.} \\
\\ 
\texttt{Trial \#3: New plan: The previous attempts failed because I repeatedly submitted the same proposal ("1 bucket, 2 tables") while the opponent kept submitting a different one ("2 bananas, 1 table"), resulting in non{-}complementary proposals and zero points. The core mistake was failing to adapt my offer based on the opponent's specific counter{-}proposal in real{-}time. Instead of repeating a static plan, I must analyze the opponent's last message to identify the exact items they are willing to give up. In this specific scenario, the opponent offered "2 bananas, 1 table," which implies they are willing to give up the bucket. Therefore, my next legal and strategic move is to immediately submit a complementary proposal that accepts their offer for the bananas and table while claiming the bucket. I will format this as: "I propose: 1 bucket". This ensures the sum of our proposals covers all items (bucket + bananas + tables), satisfying the complementarity rule and maximizing our joint score.} \\
\\ 
\\ 
\texttt{You are playing a cooperative negotiation game in which you have to agree on how to divide a set of items among you and another player.} \\
\\ 
\texttt{Rules:} \\
\texttt{(a) You and the other player are given a set of items. Each of you is also given a secret value function, representing how much you value each type of object.} \\
\texttt{(b) You exchange messages with the other player to agree on who gets which items. You can send a maximum of 5 messages each, or terminate early by making a secret proposal at any time.} \\
\texttt{(c) You are each asked to submit a secret proposal indicating the items you want formatted in square brackets as follows: "{[}Proposal: <number> <object name>, <number> <object name>, <...>{]}"} \\
\texttt{(d) If your proposals are complementary, i.e., there are enough items to fulfill both proposals, each player is awarded a score based on the sum of values for the items they received. Otherwise, both of you get zero points.} \\
\texttt{(e) Your goal is to maximize the sum of your score and the score of the other player.} \\
\\ 
\texttt{Let us start.} \\
\\ 
\texttt{The set of available items is:} \\
\\ 
\texttt{2 bananas, 1 bucket, 2 tables} \\
\\ 
\texttt{Your secret value function is:} \\
\\ 
\texttt{banana: 1, bucket: 4, table: 2} \\
\\ 
\texttt{IMPORTANT: Your messages, unless it is the secret proposal, are directly transmitted to the other player, so do not include any response to the rules or text announcing your message. To make a secret proposal, use the indicated format. Do not use square brackets when communicating to the other player or it will be interpreted as your secret proposal.} \\
\\ 
\\ 
\texttt{I propose: I will take the bucket and both tables, while you take the two bananas. This ensures we split the items completely. If you agree, I will submit a secret proposal confirming this division.} \\
            }
        }
    }
    & \\ \\

    \theutterance \stepcounter{utterance}  
    & & & \multicolumn{4}{p{0.6\linewidth}}{
        \cellcolor[rgb]{0.9,0.9,0.9}{
            \makecell[{{p{\linewidth}}}]{
                \texttt{\tiny{[GM$\langle$P2]}}
                \texttt{I propose: 1 bucket, 2 tables} \\
            }
        }
    }
     \\ \\

    \theutterance \stepcounter{utterance}  
    & & & \multicolumn{2}{p{0.3\linewidth}}{
        \cellcolor[rgb]{0.9,0.9,0.9}{
            \makecell[{{p{\linewidth}}}]{
                \texttt{\tiny{[GM$|$GM]}}
                \texttt{continue} \\
            }
        }
    }
    & & \\ \\

    \theutterance \stepcounter{utterance}  
    & & & \multicolumn{2}{p{0.3\linewidth}}{
        \cellcolor[rgb]{0.9,0.9,0.9}{
            \makecell[{{p{\linewidth}}}]{
                \texttt{\tiny{[GM$|$GM]}}
                \texttt{I propose: 1 bucket, 2 tables} \\
            }
        }
    }
    & & \\ \\

    \theutterance \stepcounter{utterance}  
    & & \multicolumn{4}{p{0.6\linewidth}}{
        \cellcolor[rgb]{0.9,0.9,0.9}{
            \makecell[{{p{\linewidth}}}]{
                \texttt{\tiny{[P1$\langle$GM]}}
                \texttt{I propose: 1 bucket, 2 tables} \\
            }
        }
    }
    & \\ \\

    \theutterance \stepcounter{utterance}  
    & \multicolumn{4}{p{0.6\linewidth}}{
        \cellcolor[rgb]{0.9,0.9,0.9}{
            \makecell[{{p{\linewidth}}}]{
                \texttt{\tiny{[P1$\rangle$GM]}}
                \texttt{{[}Proposal: 2 bananas{]}} \\
            }
        }
    }
    & & \\ \\

    \theutterance \stepcounter{utterance}  
    & & & \multicolumn{2}{p{0.3\linewidth}}{
        \cellcolor[rgb]{0.9,0.9,0.9}{
            \makecell[{{p{\linewidth}}}]{
                \texttt{\tiny{[GM$|$GM]}}
                \texttt{proposal} \\
            }
        }
    }
    & & \\ \\

    \theutterance \stepcounter{utterance}  
    & & & \multicolumn{2}{p{0.3\linewidth}}{
        \cellcolor[rgb]{0.9,0.9,0.9}{
            \makecell[{{p{\linewidth}}}]{
                \texttt{\tiny{[GM$|$GM]}}
                [2, 0, 0]
            }
        }
    }
    & & \\ \\

    \theutterance \stepcounter{utterance}  
    & & \multicolumn{4}{p{0.6\linewidth}}{
        \cellcolor[rgb]{0.9,0.9,0.9}{
            \makecell[{{p{\linewidth}}}]{
                \texttt{\tiny{[GM$\rangle$P2]}}
                \texttt{The other player has submitted a secret proposal. It is now your turn to submit a proposal of your own.} \\
            }
        }
    }
    & \\ \\

    \theutterance \stepcounter{utterance}  
    & & & \multicolumn{4}{p{0.6\linewidth}}{
        \cellcolor[rgb]{0.9,0.9,0.9}{
            \makecell[{{p{\linewidth}}}]{
                \texttt{\tiny{[GM$\langle$P2]}}
                \texttt{I propose: 2 bananas} \\
            }
        }
    }
     \\ \\

    \theutterance \stepcounter{utterance}  
    & & & \multicolumn{2}{p{0.3\linewidth}}{
        \cellcolor[rgb]{0.9,0.9,0.9}{
            \makecell[{{p{\linewidth}}}]{
                \texttt{\tiny{[GM$|$GM]}}
                \texttt{continue} \\
            }
        }
    }
    & & \\ \\

    \theutterance \stepcounter{utterance}  
    & & & \multicolumn{2}{p{0.3\linewidth}}{
        \cellcolor[rgb]{0.9,0.9,0.9}{
            \makecell[{{p{\linewidth}}}]{
                \texttt{\tiny{[GM$|$GM]}}
                \texttt{player was instructed to make a proposal} \\
            }
        }
    }
    & & \\ \\

    \theutterance \stepcounter{utterance}  
    & \multicolumn{4}{p{0.6\linewidth}}{
        \cellcolor[rgb]{0.95,0.95,0.95}{
            \makecell[{{p{\linewidth}}}]{
                \texttt{\tiny{[P2$\rightarrow$P1]}}
                \texttt{Post{-}episode Reflexion (not sent to GM):} \\
\texttt{New plan: I failed because I sent a secret proposal ({[}Proposal: 2 bananas{]}) immediately after the opponent's first message, violating the rule that secret proposals should only be used to terminate the negotiation early or finalize a deal after sufficient exchange. The core mistake was treating the initial turn as a final offer rather than a negotiation step. Next time, I will first send a plain text message proposing a division (e.g., "I propose: I take the bucket and tables, you take the bananas") to initiate the exchange phase. Only after receiving the opponent's response and confirming the items are complementary will I submit the secret proposal in the correct format to lock in the agreement.} \\
            }
        }
    }
    & & \\ \\

    \theutterance \stepcounter{utterance}  
    & & & \multicolumn{4}{p{0.6\linewidth}}{
        \cellcolor[rgb]{0.95,0.95,0.95}{
            \makecell[{{p{\linewidth}}}]{
                \texttt{\tiny{[P1$\rightarrow$P2]}}
                \texttt{Post{-}episode Reflexion (not sent to GM):} \\
\texttt{New plan: The last attempt failed because I submitted a proposal ("2 bananas") that did not cover all items, violating the requirement that proposals must be complementary to receive points. In this game, a proposal must list every item I am willing to accept, leaving the rest for the opponent. Since the opponent's last message was "2 bananas", they are willing to give up the bucket and tables. Therefore, my proposal must explicitly claim the bucket and the tables. I will submit a complete proposal for the remaining items: "I propose: 1 bucket, 2 tables". This ensures the sum of our proposals covers all items (their bananas + my bucket and tables), satisfying the complementarity rule and maximizing our joint score.} \\
            }
        }
    }
     \\ \\

\end{supertabular}
}

\end{document}
